\section{模型建立}

\subsection{问题描述}

本文研究动态需求下多家配送企业的协同配送问题。各企业拥有独立车场和油电混合车队，
车辆从所属车场出发并返回所属车场，可直接服务其他企业的客户，但不在车场之间转运货物；
同一车辆可在一个运营日内连续执行多趟配送任务；电动车可在车场或公共充电站部分充电，
充电功率随电池荷电状态变化，充电费用和间接碳排放分别由实际充电时段的电价和电网碳强度确定。
动态订单到达后，系统在批事件处理时刻对尚未完成的任务重新规划。
模型的优化目标是在满足载重、时间窗、电量和趟次衔接等约束的前提下，
合理确定车辆路径、车型选择和充电安排，使配送总成本最小。

模型假设条件如下：1）客户、车场、车辆及已到达订单的信息已知；
2）每个客户仅由一辆车服务一次，需求不可拆分；
3）车辆载重和电量不得超过相应容量限制；
4）同一车辆相邻两趟在时间上不得重叠；
5）客户具有硬时间窗，车辆提前到达时允许等待；
6）分时电价和时变电网碳强度为外生参数；
7）协同路径确定后再进行联盟成本分摊，不以分摊结果改变路径模型。
模型符号说明如表\ref{tab:symbols}所示。

\begin{table}[H]
\centering
\caption{符号说明}
\label{tab:symbols}
\fontsize{8.3pt}{9.3pt}\selectfont
\renewcommand{\arraystretch}{0.88}
\setlength{\tabcolsep}{1.5pt}
\begin{tabular*}{\linewidth}{@{\extracolsep{\fill}}cJ|cJ@{}}
\toprule
符号 & 说明 & 符号 & 说明\\
\midrule
$H$ & 参与协作的企业集合 & $D$ & 车场集合\\
$N$ & 客户集合 & $R$ & 充电站集合\\
$\widehat R$ & 充电节点副本集合 & $D_{N+1}$ & 终点车场副本集合\\
$0$ & 车辆所属车场的起终点 & $O$ & 车场与客户节点集合\\
$A$ & 车场、客户与充电节点集合 & $K$ & 车辆集合\\
$K^C$ & 燃油车集合 & $K^E$ & 电动车集合\\
$W$ & 趟次集合 & $\mathcal M$ & 分时核算时段集合\\
$q_i$ & 客户$i$的需求量 & $[e_i,l_i]$ & 客户$i$的时间窗\\
$s_i$ & 客户$i$的服务时间 & $d_{ij}$ & 节点$i$与$j$之间的距离\\
$t_{ij}$ & 节点$i$与$j$之间的行驶时间 & $Q$ & 车辆最大载重\\
$B$ & 电池容量 & $T^H$ & 车辆运营时域\\
$M$ & 足够大的正数 & $t_i^R$ & 客户$i$的需求释放时刻\\
$t_k^D$ & 车辆$k$所属车场的作业时间 & $t_{k,i}$ & 车场与客户$i$之间的行驶时间\\
$x_{ijk}^C$ & 燃油车$k$是否行驶弧$(i,j)$ & $x_{ijk}^E$ & 电动车$k$是否行驶弧$(i,j)$\\
$\gamma_{k,w}$ & 车辆$k$是否执行第$w$趟 & $\lambda_{k,i,w}$ & 客户$i$是否由车辆$k$在第$w$趟服务\\
$\xi_{i,j,k,w}$ & 第$w$趟是否连续访问$i$和$j$ & $\mu_{k,i,w}$ & 客户$i$是否为第$w$趟首客户\\
$\eta_{k,i,w}$ & 客户$i$是否为第$w$趟末客户 & $\tau_{k,i,w}$ & 车辆$k$第$w$趟到达客户$i$的时刻\\
$\tau_{k,w}^{L}$ & 第$w$趟的出发时刻 & $\tau_{k,w}^{B}$ & 第$w$趟的返回时刻\\
$\delta_{k,w}^{L}$ & 第$w$趟的出发载重 & $L_{ij}$ & 燃油车弧段油耗量\\
$E_{ij}$ & 燃油车弧段碳排放量 & $c^C$ & 燃油车启动成本\\
$c^E$ & 电动车启动成本 & $m^C$ & 燃油车单位里程成本\\
$m^E$ & 电动车单位里程成本 & $w_r$ & 单位充电量成本\\
$w_o$ & 单位燃油成本 & $w_c$ & 单位排放成本\\
$c$ & 单位碳价 & $P_m^{\mathrm{TOU}}$ & 时段$m$的充电电价\\
$b_{jm}$ & 充电节点$j$是否在时段$m$被访问 & $rc_{jm}$ & 节点$j$在时段$m$的充电量\\
$y_j^A$ & 到达充电节点$j$时的电量 & $y_j^D$ & 离开充电节点$j$时的电量\\
$z_j$ & 到达充电节点$j$的时刻 & $r_i$ & 电动车在充电节点$i$的充电量\\
$N'$ & 更新后的待配送客户集合 & $K^{C'}$ & 更新后可用燃油车集合\\
$K^{E'}$ & 更新后可用电动车集合 & $K^{\mathrm{on}}$ & 已出发且可继续执行任务的车辆集合\\
$Z^C$ & 重新派出的燃油车集合 & $Z^E$ & 重新派出的电动车集合\\
$u(t)$ & 第$t$个批事件处理时刻 & $u_q$ & 定量触发时刻\\
$T$ & 动态需求处理间隔 & $\bar q$ & 累计需求上限\\
$F_C$ & 燃油车排放总量 & $F_E$ & 电动车排放总量\\
$C(S)$ & 联盟$S$的配送成本 & $\psi_h$ & 企业$h$的成本分摊额\\
$S$ & 企业子联盟 & $h$ & 企业索引\\
\bottomrule
\end{tabular*}
\end{table}

\Needspace{10\baselineskip}
\subsection{目标函数}

\subsubsection{车辆启动成本 $F_1$}

车辆启动成本与投入运营的车辆数有关\cite{ref:71}：
\begin{equation}
F_1=c^C\sum_{k\in K^C}\sum_{j\in O}x_{0jk}^C
+c^E\sum_{k\in K^E}\sum_{j\in A}x_{0jk}^E.
\label{eq:qiu-f1}
\end{equation}

\subsubsection{行驶成本 $F_2$}

行驶成本与车辆行驶里程成正比\cite{ref:71}：
\begin{equation}
F_2=m^C\sum_{k\in K^C}\sum_{(i,j)\in O}x_{ijk}^Cd_{ij}
+m^E\sum_{k\in K^E}\sum_{(i,j)\in A}x_{ijk}^Ed_{ij}.
\label{eq:qiu-f2}
\end{equation}

\subsubsection{充电成本 $F_3$}

充电成本由各时段的充电量和对应电价确定\cite{ref:96}：
\begin{equation}
F_3=\sum_{m\in \mathcal M}\sum_{j\in\widehat R\cup D_{N+1}}
P_m^{\mathrm{TOU}}rc_{jm}.
\label{eq:deng-41}
\end{equation}

\subsubsection{油耗成本 $F_4$}

油耗成本与燃油车的弧段油耗成正比\cite{ref:71}：
\begin{equation}
F_4=w_o\sum_{k\in K^C}\sum_{(i,j)\in O}x_{ijk}^CL_{ij}.
\label{eq:qiu-f4}
\end{equation}

\subsubsection{碳排放成本 $F_5$}

碳排放成本由燃油车直接排放和电动车充电间接排放共同确定\cite{ref:95}：
\begin{equation}
F_5=c(F_C+F_E).
\label{eq:gong-carbon-cost}
\end{equation}

\subsection{混合车队能耗模型}
\label{sec:energy-model}

采用实际能耗模型计算车辆行驶产生的电耗和油耗\cite{ref:23}。车辆$k$在弧$(i,j)$上的机械功率为
\begin{equation}
T_{ij}^{k}(u_{ij}^{k},v_{ij}^{k})
=\frac{1}{2}c_d\rho^a A_k^f(v_{ij}^{k})^3
+(m_k+u_{ij}^{k})gc_rv_{ij}^{k}.
\label{eq:chen-power}
\end{equation}
其中，$u_{ij}^{k}$和$v_{ij}^{k}$分别为弧段载重和速度，$c_d,\rho^a,A_k^f,g,c_r,m_k$
依次为空气阻力系数、空气密度、迎风面积、重力加速度、滚动阻力系数和空车质量。
电动车的弧段电耗为
\begin{equation}
b_{ij}^{k}(u_{ij}^{k},v_{ij}^{k})
=\phi^{d}\varphi^{d}T_{ij}^{k}(u_{ij}^{k},v_{ij}^{k})\frac{d_{ij}}{v_{ij}^{k}},
\label{eq:chen-electricity}
\end{equation}
其中，$\phi^d$和$\varphi^d$分别为电动机效率和电池放电效率。燃油车的单位时间油耗率为
\begin{equation}
G_{ij}^{k}(u_{ij}^{k},v_{ij}^{k})
=\xi^{\mathrm f}\frac{\varepsilon R^{\mathrm e}+T_{ij}^{k}(u_{ij}^{k},v_{ij}^{k})/\eta}{\kappa\psi},
\label{eq:chen-fuel-rate}
\end{equation}
其中，$\xi^{\mathrm f},\varepsilon,R^{\mathrm e},\eta,\kappa,\psi$分别为柴油空气质量比、
摩擦因数、发动机参数、发动机效率、柴油热值和燃油率换算系数。弧段油耗为
\begin{equation}
f_{ij}^{k}(u_{ij}^{k},v_{ij}^{k})
=G_{ij}^{k}(u_{ij}^{k},v_{ij}^{k})\frac{d_{ij}}{v_{ij}^{k}}.
\label{eq:chen-fuel}
\end{equation}

\subsection{非线性充电函数}
\label{sec:charging-model}

设$BC_i=\{0,\ldots,bc_i\}$为充电函数断点集合，$q_{i,k}$和$\epsilon_{i,k}$
分别为断点电量和累计充电时间，$\rho_{i,k}$为相邻断点间的充电斜率。
电量与充电时间的关系采用Wang等\cite{ref:94}式（A.1）--（A.2）：
\begin{equation}
\phi_i(t)=
\begin{cases}
\rho_{i,0}t, & t\le\epsilon_{i,0},\\
q_{i,0}+\rho_{i,1}(t-\epsilon_{i,0}), & \epsilon_{i,0}<t\le\epsilon_{i,1},\\
\cdots & \cdots\\
q_{i,k-1}+\rho_{i,k}(t-\epsilon_{i,k-1}),
& \epsilon_{i,k-1}<t\le\epsilon_{i,k},\\
\cdots & \cdots
\end{cases}
\label{eq:wang-charge-function}
\end{equation}
\begin{equation}
\phi_i(q,\Delta t_1,\Delta t_2)
=\phi_i\!\left(\phi_i^{-1}(q)+\Delta t_2\right)
-\phi_i\!\left(\phi_i^{-1}(q)+\Delta t_1\right),
\quad
\phi_i^{-1}(q)\le\Delta t_1\le\Delta t_2\le\epsilon_{i,bc_i}.
\label{eq:charge_duration}
\end{equation}
其中，$q$为开始充电时的电量，$\Delta t_1$和$\Delta t_2$为两个充电时长；
式（\ref{eq:charge_duration}）表示两充电时长之间的充电量差。

\subsection{分时充电与碳排放核算}
\label{sec:time-carbon-model}

设$[\underline T_m,\overline T_m]$为时段$m$的起止时刻，$b_{jm}$表示充电节点$j$
是否在时段$m$被访问，$rc_{jm}$表示相应充电量。时段和充电量采用Deng等\cite{ref:96}
式（35）--（39）确定。$V_0$和$\widehat V$为相应前驱节点集合，
$x_{ij}$为充电层弧选择变量：
\begin{align}
\sum_{m\in \mathcal M}\underline T_m b_{jm}<z_j
&\le\sum_{m\in \mathcal M}\overline T_m b_{jm},
&&\forall j\in \widehat R\cup D_{N+1},
\label{eq:deng-35}\\
\sum_{m\in \mathcal M}b_{jm}&=\sum_{i\in V_0}x_{ij},
&&\forall j\in\widehat R,
\label{eq:deng-36}\\
\sum_{m\in \mathcal M}b_{jm}&=\sum_{i\in\widehat V}x_{ij},
&&\forall j\in D_{N+1},
\label{eq:deng-37}\\
0\le rc_{jm}&\le B b_{jm},
&&\forall j\in\widehat R\cup D_{N+1},\ m\in \mathcal M,
\label{eq:deng-38}\\
\sum_{m\in \mathcal M}rc_{jm}&=y_j^D-y_j^A,
&&\forall j\in\widehat R\cup D_{N+1}.
\label{eq:deng-39}
\end{align}

燃油车直接排放采用巩亮等\cite{ref:95}的碳排放因子法：
\begin{align}
F_j&=\frac{12}{44}Q_jU_jO_j,
\label{eq:gong-factor}\\
C_{i1}&=\sum_j K_{i,j}\times F_j,
\label{eq:gong-direct}\\
F_C&=\delta C_{i1}.
\label{eq:gong-fc}
\end{align}
其中，$F_j$为$j$类化石燃料的碳排放因子，$Q_j,U_j,O_j$分别为低位发热量、
单位热值含碳量和碳氧化率，$K_{i,j}$为第$i$种交通方式对$j$类燃料的消耗量，
$\delta$为碳排放因子，单位为$\mathrm{kgCO_2/L}$。电动车充电的间接排放为\cite{ref:96}
\begin{align}
IC_m(d)&=CO_{2eq}^{\mathrm{Wind}}Pt_m^{\mathrm{Wind}}(d)
+CO_{2eq}^{\mathrm{Others}}Pt_m^{\mathrm{Others}}(d),
\label{eq:deng-intensity}\\
CO_{2eq}(d)&=\sum_{m\in \mathcal M}IC_m(d)r_m(d).
\label{eq:deng-emission}
\end{align}
其中，$IC_m(d)$为第$d$日时段$m$的电网碳强度，$r_m(d)$为相应时段的充电量；
$CO_{2eq}(d)$记为电动车排放总量$F_E$。

\subsection{动态信息处理}
\label{sec:dynamic-demand-model}

本文考虑客户需求新增、取消和改量三类动态信息，其中改量先取消原订单，再按新增订单处理。
设$k'$为重规划车辆索引。动态阶段的五项成本为\cite{ref:71}
\begin{align}
F_1'&=F_1+c^C\sum_{k\in Z^C}\sum_{j\in N'}x_{0jk}^C
+c^E\sum_{k\in Z^E}\sum_{j\in N'}x_{0jk}^E,
\label{eq:qiu-dyn-f1}\\
F_2'&=m^C\sum_{k'\in K^{\mathrm{on}}\cup K^{C'}}\sum_{(i,j)\in N\cup N'}x_{ijk'}^Cd_{ij}
+m^E\sum_{k'\in K^{\mathrm{on}}\cup K^{E'}}\sum_{(i,j)\in A\cup N'}x_{ijk'}^Ed_{ij},
\label{eq:qiu-dyn-f2}\\
F_3'&=w_r\sum_{k'\in K^{\mathrm{on}}\cup K^{E'}}\sum_{(i,j)\in N'\cup A}x_{ijk'}^Er_i,
\label{eq:qiu-dyn-f3}\\
F_4'&=w_o\sum_{k'\in K^{\mathrm{on}}\cup K^{C'}}\sum_{(i,j)\in O\cup N'}x_{ijk'}^CL_{ij},
\label{eq:qiu-dyn-f4}\\
F_5'&=w_c\sum_{k'\in K^{\mathrm{on}}\cup K^{C'}}\sum_{(i,j)\in O\cup N'}x_{ijk'}^CE_{ij},
\label{eq:qiu-dyn-f5}\\
\min F'&=F_1'+F_2'+F_3'+F_4'+F_5'.
\label{eq:qiu-dyn-obj}
\end{align}
设$[t_0,t_n]$为接收动态需求的时间区间，$T$为处理间隔，$\bar q$为累计需求上限。
定时与定量相结合的批事件处理时刻为
\begin{align}
u(t)&=\min\{u_q,u(t-1)+T,t_n\},
\label{eq:qiu-trigger-time}\\
u_q&=\min\{u_q\mid q(u(t-1),u_q)\ge\bar q,\ u(t-1)>u_q\}.
\label{eq:qiu-trigger-quantity}
\end{align}
其中，$q(a,b)$表示区间$[a,b]$内到达的累计需求，$u(0)=t_0$，$u(n)=t_n$。
在批事件处理时刻，已完成任务不再进入下一次规划，
尚未完成的客户和可用车辆重新求解。

\subsection{协作配送成本分摊}
\label{sec:allocation-model}

协作路径确定后，采用Shapley值分配联盟成本\cite{ref:91}。设$S\subseteq H$为企业子联盟，
$C(S)$为联盟$S$的配送成本，则企业$h$的成本分摊额为
\begin{equation}
\psi_h=\sum_{S\subseteq H\setminus\{h\}}
\frac{|S|!(|H|-|S|-1)!}{|H|!}
\bigl[C(S\cup\{h\})-C(S)\bigr].
\label{eq:shapley-cost}
\end{equation}
式（\ref{eq:shapley-cost}）按企业加入联盟的边际成本确定分摊额，且满足
$\sum_{h\in H}\psi_h=C(H)$。

\subsection{模型建立}

综上，模型的目标函数为
\begin{equation}
\min F=F_1+F_2+F_3+F_4+F_5.
\label{eq:obj}
\end{equation}
车辆从所属车场出发并返回所属车场，同一车辆可连续执行多趟任务。
多车场多趟约束采用Zhen等\cite{ref:44}式（2）--（20）：
\begin{align}
\sum_{k\in K}\sum_{w\in W}\lambda_{k,i,w}&=1,
&&\forall i\in N,
\label{eq:zhen-2}\\
\mu_{k,l,w}+\sum_{i\in N}\xi_{i,l,k,w}
&=\eta_{k,l,w}+\sum_{j\in N}\xi_{l,j,k,w}
=\lambda_{k,l,w},
&&\forall l\in N,k\in K,w\in W,
\label{eq:zhen-3}\\
\sum_{i\in N}\mu_{k,i,w}
&=\sum_{i\in N}\eta_{k,i,w}=\gamma_{k,w},
&&\forall k\in K,w\in W,
\label{eq:zhen-4}\\
\gamma_{k,w+1}&\le\gamma_{k,w},
&&\forall k\in K,w\in W\setminus\{|W|\},
\label{eq:zhen-5}\\
\sum_{i\in N}\lambda_{k,i,w}&\le\gamma_{k,w}|N|,
&&\forall k\in K,w\in W,
\label{eq:zhen-6}\\
\delta_{k,w}^{L}&=\sum_{i\in N}\lambda_{k,i,w}q_i,
&&\forall k\in K,w\in W,
\label{eq:zhen-7}\\
\delta_{k,w}^{L}&\le\gamma_{k,w}Q,
&&\forall k\in K,w\in W,
\label{eq:zhen-8}\\
\tau_{k,i,w}&\le\lambda_{k,i,w}l_i,
&&\forall i\in N,k\in K,w\in W.
\label{eq:zhen-9}
\end{align}
\begin{align}
\tau_{k,w}^{L}&\ge\lambda_{k,i,w}(t_i^{R}+t_k^{D})-M(1-\gamma_{k,w}),
&&\forall i\in N,k\in K,w\in W,
\label{eq:zhen-10}\\
\tau_{k,w+1}^{L}&\ge\tau_{k,w}^{B}+t_k^{D}-M(1-\gamma_{k,w+1}),
&&\forall i\in N,k\in K,w\in W,
\label{eq:zhen-11}\\
\tau_{k,i,w}&\ge\tau_{k,w}^{L}+t_{k,i}-M(1-\mu_{k,i,w}),
&&\forall i\in N,k\in K,w\in W,
\label{eq:zhen-12}\\
\tau_{k,w}^{B}&\ge\max\{\tau_{k,i,w},e_i\}+s_i+t_{k,i}-M(1-\eta_{k,i,w}),
&&\forall i\in N,k\in K,w\in W,
\label{eq:zhen-13}\\
\tau_{k,j,w}&\ge\max\{\tau_{k,i,w},e_i\}+s_i+t_{ij}-M(1-\xi_{i,j,k,w}),
&&\forall i,j\in N,k\in K,w\in W,
\label{eq:zhen-14}\\
\tau_{k,w}^{B}&\le T^{H},
&&\forall k\in K,w\in W.
\label{eq:zhen-15}
\end{align}
\begin{align}
\gamma_{k,w}&\in\{0,1\},
&&\forall k\in K,w\in W,
\label{eq:zhen-16}\\
\xi_{i,j,k,w}&\in\{0,1\},
&&\forall i,j\in N,k\in K,w\in W,
\label{eq:zhen-17}\\
\lambda_{k,i,w}&\in\{0,1\},
&&\forall i\in N,k\in K,w\in W,
\label{eq:zhen-18}\\
\mu_{k,i,w}&\in\{0,1\},
&&\forall i\in N,k\in K,w\in W,
\label{eq:zhen-19}\\
\eta_{k,i,w}&\in\{0,1\},
&&\forall i\in N,k\in K,w\in W.
\label{eq:zhen-20}
\end{align}
式（\ref{eq:zhen-2}）保证每个客户仅被服务一次；式（\ref{eq:zhen-3}）表示趟内客户访问平衡；
式（\ref{eq:zhen-4}）限定每个已执行趟次仅有一个首客户和一个末客户；
式（\ref{eq:zhen-5}）保证车辆按趟次顺序连续执行任务；式（\ref{eq:zhen-6}）限定客户只能分配给已启用趟次；
式（\ref{eq:zhen-7}）计算车辆每趟的出发载重，式（\ref{eq:zhen-8}）为车辆载重约束；
式（\ref{eq:zhen-9}）为客户时间窗约束；式（\ref{eq:zhen-10}）表示车辆各趟出发时刻与客户需求释放时刻及车场作业时间之间的关系；
式（\ref{eq:zhen-11}）表示同一车辆相邻趟次的时间衔接；
式（\ref{eq:zhen-12}）和式（\ref{eq:zhen-13}）分别表示车辆到达首客户和返回车场的时间关系；
式（\ref{eq:zhen-14}）表示趟内相邻客户的时间递推；式（\ref{eq:zhen-15}）限制车辆返回时刻；
式（\ref{eq:zhen-16}）--（\ref{eq:zhen-20}）规定决策变量的取值范围。
车辆所属车场决定$t_{k,i}$和$t_k^D$，车辆可直接服务其他车场的客户，模型中不设置车场间货物转运变量。
